\documentclass[../../main.tex]{subfiles}
% \documentclass[a4paper,10pt]{ltjsarticle}
\begin{document}
\rhead{京都大学後期$1997$年$1$番}
\setcounter{equation}{0}
\begin{oframed}
  \documentclass[../../main.tex]{subfiles}
\begin{document}

$2$つの放物線$C_1:y=x^2$，$C:y=ax^2(a\neq 1)$を考える．
$C$を$x$軸方向に$p$，ついで$y$軸方向に$q$だけ平行移動した放物線を$C_{p,q}$と表す．
次の条件（*）を満たすような$p$，$q$が存在するための$a$の範囲を求めよ．
\begin{itemize}
  \item[（*）] $C_{p,q}$は$C_1$と$2$点で交わる．
    $1$つの交点は$C_{p,q}$の頂点であり，他の交点においては両者の接線は直交する．
\end{itemize}

\end{document}

\end{oframed}
\setlength{\columnseprule}{0.4pt}
\begin{multicols}{2}
  {\bf [解答]}

  \begin{figure}[H]
    \centering
    \begin{tikzpicture}[x=1.8cm,y=1.35cm,>=stealth,
      every node/.style={font=\small}]
      % a=-1/2, p=3/(2sqrt(2)), q=p^2=9/8 の場合．
      \draw[->,thick,gray!80!black] (-1.25,0) -- (1.8,0) node[right] {$x$};
      \draw[->,thick,gray!80!black] (0,-0.8) -- (0,2.4) node[above] {$y$};
      \draw[very thick,blue!85!black,domain=-1.1:1.5,samples=80]
      plot (\x,{\x*\x});
      \draw[very thick,teal!85!black,domain=-1.1:2,samples=80]
      plot (\x,{-0.5*(\x-1.06066)*(\x-1.06066)+1.125});
      \coordinate (V) at (1.06066,1.125);
      \coordinate (T) at (-0.35355,0.125);
      \draw[densely dashed,thick,blue!85!black] (-1,0.582) -- (0.35,-0.372);
      \draw[densely dashed,thick,teal!85!black] (-1,-0.789) -- (0.35,1.12);
      \node[circle,fill=red!85!black,inner sep=1.3pt] at (V) {};
      \node[circle,fill=red!85!black,inner sep=1.3pt] at (T) {};
      \node[above left=2pt] at (V) {頂点$(p,q)$};
      \node[above left=2pt] at (T) {直交する接線};
      \node[blue!85!black,right] at (0.9,2) {$C_1$};
      \node[teal!85!black,right] at (1.65,0.95) {$C_{p,q}$};
    \end{tikzpicture}
    \caption{$C_1$と$C_{p,q}$の2交点と接線の直交}
    \label{1997-1:fig:1}
  \end{figure}

  題意より，$C_{p,q}$の方程式は
  \begin{align}
    & C_{p,q}: y=a(x-p)^2+q & (a \neq 1) \label{eq:1}
  \end{align}
  である．ここで，$p0$とすると，$C_{p,q}$と$C_{1}$を連立して
  \begin{align}
    ax^2 + q = x^2
    \quad\therefore
    x =
    \begin{cases}
      \pm \sqrt{\dfrac{q}{a-1}} & (q(a-1)\ge 0) \\
      \text{解なし} & (q(a-1)<0)
    \end{cases}
  \end{align}
  となる．これが$C_{p=0,q}$の頂点$(0,q)$であるには$q(a-1) \ge 0$の元で$\sqrt{\dfrac{q}{a-1}}=0$つまり$q=0$であることが必要だが，この時は$C_{p,q}$と$C_{1}$は一点でしか交わらないため不適である．従って以下
  \begin{align}
    p \neq 0
  \end{align}
  として考える．

  $C_{p,q}$の頂点は$(p,q)$である．
  この点で$C_1$および$C_{p,q}$が交わることから，
  $C_1$は$(p,q)$を通る．従って
  \begin{align}
    q & = p^2\label{eq:2}
  \end{align}
  である．
  この元で，$C_{p,q}$と$C_{1}$を連立した$x$の二次方程式
  \begin{align}
    x^2 = a(x-p)^2+q
  \end{align}
  のもう一つの解の点で$C_1$,$C_{p,q}$の接線が直交する．
  \cref{eq:2}を代入すると
  \begin{align}
    & (a-1)x^2-2apx+(a+1)p^2 = 0                               \\
    & (a-1)\left(x-p\right)\left(x-\dfrac{a+1}{a-1}p\right) = 0
  \end{align}
  となる．$a-1=0$の時は方程式が一次方程式となり，二点で交わる条件を満たさないので不適である．
  すなわち
  \begin{align}
    a \neq 1 \label{eq:3}
  \end{align}
  である．
  この元で，もう一つの解は
  \begin{align}
    x = \dfrac{a+1}{a-1}p
  \end{align}
  である．この点で$C_1$, $C_{p,q}$の接線の傾きは
  \begin{align}
    & C_1: 2\dfrac{a+1}{a-1}p                     \\
    & C_{p,q}: 2a\left(\dfrac{a+1}{a-1}p-p\right)
  \end{align}
  だから，これらが直交する条件より，
  \begin{align}
    & 2a\left(\dfrac{a+1}{a-1}p-p\right) \cdot 2\dfrac{a+1}{a-1}p = -1 \\
    & 8a \dfrac{a+1}{(a-1)^2} p^2 = -1
  \end{align}
  $a=0, a=-1$はこの条件を満たさないので不適で，以下
  \begin{align}
    \begin{dcases}
      & a \neq 0  \\
      & a \neq -1
    \end{dcases}\label{eq:4}
  \end{align}
  として考えると，この方程式を$p$について解いて
  \begin{align}
    p^2 = -\dfrac{(a-1)^2}{8a(a+1)}
  \end{align}
  このような実数$p\neq 0$が存在するには，右辺が$0$より大きければよく，
  \begin{align}
    -\dfrac{(a-1)^2}{8a(a+1)} > 0
  \end{align}
  が求める$a$の条件である．
  \cref{eq:3}より$a\neq 1$だから分子は負であり，よってこの条件を満たすには
  分母が負であれば良い．
  \begin{align}
    & a(a+1) < 0 \\
    \therefore
    & -1 < a < 0
  \end{align}
  これは\cref{eq:3}の条件を満たしているから，求める条件は
  \begin{align}
    -1 < a < 0
  \end{align}
  である． $\cdots$(答)

  \vspace{10pt}
  {\bf [解説]}

  \newpage
\end{multicols}
\end{document}
